\svnidlong{$LastChangedBy$}{$LastChangedRevision$}{$LastChangedDate$}{$HeadURL: http://freevariable.com/dissertation/branches/diss-template/frontmatter/abstract.tex $}
\vcinfo{}

Ultra-high energy (UHE, $E>100$~PeV) neutrinos are a rare piece of multi-messenger astrophysics and are especially valuable for understanding the UHE cosmic ray flux and composition. 
When UHE neutrinos interact in dielectric media, they initiate particle cascades that emit optical-wavelength Cherenkov and radio-wavelength Askaryan emission.
The highest energy neutrino ever detected had an energy of $\sim220$~PeV and was found by KM3NeT via Cherenkov Radiation in the Mediterranean Sea emitted by the $\sim120$~PeV muon created in the parent neutrino's interaction~\citep{km3net_neutrino}. 
The IceCube Neutrino Observatory has also observed multiple neutrinos with energies over 1~PeV via Cherekov Radiation in the Antarctic Ice Sheet at the South Pole~\citep{icecube_ehe}.
Monitoring ice for Askaryan Radiation is also promising due to long attenuation and scattering lengths for radio emission in ice, allowing a single detector to observe UHE neutrino interactions many kilometers away. 
Radio-based UHE neutrino detectors have observed Askaryan emission in ice from cosmic ray showers \citep{ara_cr, arianna_cr, anita_cr} but no radio-based detector has observed a neutrino candidate yet. 

This thesis comprises simulations and analyses dedicated to finding UHE neutrinos using in-ice detectors at the South Pole and Greenland. 
Chapter~\ref{chap:intro} introduction to the field of UHE neutrino observations and it's place in multi-messenger astrophysics. 
A simulation study is then presented in Chapter~\ref{chap:gen2_infill} quantifying the benefits of joint Askaryan and Cherenkov emission observations in the proposed expansion of the IceCube Neutrino Observatory by placing radio antennas above all planned optical modules. 
Next, Chapter~\ref{chap:ARA} describes the Askaryan Radio Array (ARA) and a literature review of previous results from the ARA collaboration is presented.
The following Chapter~\ref{chap:veff} details simulation studies of ARA's sensitivity to UHE neutrinos, featuring the simulation of outgoing particles so that events like KM3NeT's UHE neutrino are considered in ARA's performance and characterized for potential observations. 
Chapter~\ref{chap:5sa_cleaning} and Chapter~\ref{chap:5sa_analysis} of this thesis describe the data cleaning and analysis of ARA's second station in the context of an array-wide, decadal neutrino search over the full array.
The thesis then concludes with Chapter~\ref{chap:conclusion} summarizing all discussed topics and presenting next steps and suggestions for future detectors and UHE neutrino analyses, especially in the context of the Radio Neutrino Observatory of Greenland and IceCube-Gen2 Radio \citep{rnog_design, icecube_gen2}. 

\comment{Can I share the antenna centering spacers somewhere in here?}